\section{Odd kernels and pair-offset constructions}
\label{ok:chapter}\label{app:odd-kernels}

This section supplies the path constructions behind the odd three-arm
$\alpha$ theorem and the additional odd-pair offset intervals. A word
lists the labels of a path in their order. When it describes an arm,
the hub label zero is understood to precede the word. All intervals
below are intervals of integers; an interval with its lower endpoint
larger than its upper endpoint is empty. Concatenation does not
identify the last entry of one word with the first entry of the next.

The only prescribed-path existence inputs are Lemmas~\ref{common:even-path}
and~\ref{common:odd-path}. Every use of the odd-order input checks its
exception explicitly. The small-minimum three-arm constructions are
proved in Section~\ref{app:small-alpha}. All other constructions used
here are given below, including their literal initial words.

\subsection{The kernel calculus}
\label{ok:calculus}

\begin{definition}[Odd kernel]\label{ok:kernel-definition}
For integral parameters $S,U,T,V$, with $T\ge1$, a word of type
$K(S,U,T,V)$ has $2T+1$ entries and the following inventories:
\begin{align}
 \text{lower labels}&=[1,T], \nonumber\\
 \text{upper labels}&=[T+1,2T+2]\setminus\{T+V\},\label{ok:kernel-labels}\\
 \text{internal differences}
   &=[1,2T+1]\setminus\{D\},
 \qquad D=S+T-U+V, \label{ok:kernel-colors}\\
 \text{ordered endpoints}&=(E,F),\nonumber\\
 E&=2T-S+2,\qquad F=2T+2S-U-V+3. \label{ok:kernel-ends}
\end{align}
Entries alternate between the two classes and both endpoints are
upper. In particular the displayed hole, missing color, and endpoints
must lie in their specified ranges. A kernel means an actual word
with these properties, rather than merely a feasible parameter tuple.
\end{definition}

The parameters of an auxiliary kernel need not be parameters of a
spider completion. In particular $U=0$ and $V>U$ are allowed in
source kernels. The restriction $1\le V\le U$ is imposed when an
even path of half-order $U$ is used to complete or extract a kernel.

\begin{lemma}[Completion of a kernel]\label{ok:completion}
Suppose $p$ is $K(s,u,t,v)$, where $s,t,u\ge1$ and $1\le v\le u$.
Then $\Sp{2s+1,2t+1,2u+1}$ has a root-zero $\alpha$-labeling of
index $s+t+u$. Its first arm can be chosen as the maximum Walecki word
\[
 W=(N,1,N-1,2,\ldots,N-s),\qquad N=2s+2t+2u+3.
\]
Every label on the other two arms is smaller than $N-s$.
\end{lemma}
\begin{proof}
Map a lower entry $x$ of $p$ to $x+s$, and an upper entry to
$x+s+2u$; call the resulting arm $A$. Take the even path on
$[1,2u]$ from upper endpoint $u+v$ to lower endpoint $v$, append
$2u+v$, and add $s+t$ to all of its entries; call this arm $B$.
Lemma~\ref{common:even-path} applies because $1\le v\le u$.

The lower labels of the three arms are respectively
$[1,s]$, $[s+1,s+t]$, and $[s+t+1,s+t+u]$. Put $K=s+t+u$.
The upper labels of $B$ before its appended entry are
$[K+1,s+t+2u]$. The upper labels of $A$ are
\[
 [s+t+2u+1,s+2t+2u+2]\setminus\{s+t+2u+v\};
 \]
the appended entry of $B$ fills that hole. Finally $W$ supplies
$[N-s,N]$. These are disjoint consecutive blocks, with the one
explicit hole filled once, so the nonroot labels are exactly $[1,N]$.
All edges cross $K$.

The internal colors of $A$ are
$[2u+1,2t+2u+1]\setminus\{D+2u\}$, and its root edge has color
$s+2u+E=2t+2u+2$. The internal colors of $B$ are $[1,2u]$:
its appended edge has color $2u$. Its root edge has color
\[
 s+t+u+v=D+2u.
 \]
The $2s+1$ edges of $W$, including its root edge, have consecutive
colors $[2t+2u+3,N]$. The three inventories partition $[1,N]$.
\end{proof}

\begin{lemma}[Diagonal growth]\label{ok:diagonal-growth}
Let $p$ be $K(S,U,T,V)$, put $d=T-S$, and assume $d\ge0$.
For every $P\ge d+2$ one obtains
\[
 K(S+P,U+P,T+P,V+P).
\]
\end{lemma}
\begin{proof}
Take an even path $q$ on $[1,2P]$ from upper $P+d+2$ to lower
$d+2$. Its endpoints are allowed because $1\le d+2\le P$.
Add $T$ to every entry of $q$. In the old word keep lower entries
fixed and raise upper entries by $2P$. Prepend the translated $q$.

The new lower labels are $[1,T]\cup[T+1,T+P]$.
The new upper blocks are $[T+P+1,T+2P]$ and
$[T+2P+1,2T+2P+2]\setminus\{T+V+2P\}$.
Thus the new hole is $(T+P)+(V+P)$. The prefix contributes
$[1,2P-1]$ and its last label is $T+d+2=E$.
The old first label became $E+2P$, so the join has color $2P$.
Every old color increases by $2P$. The new palette is therefore
$[1,2(T+P)+1]\setminus\{D+2P\}$, as required.
The new endpoints are
$T+P+d+2=2(T+P)-(S+P)+2$ and $F+2P$.
Substitution in \eqref{ok:kernel-ends} proves the asserted parameters.
\end{proof}

\begin{lemma}[Carrier prefix]\label{ok:carrier-prefix}
From $K(S,U,T,V)$ and any $w\ge S+1$ one obtains
\[
 K(S,U+w,T+w,V).
\]
\end{lemma}
\begin{proof}
Take the even path on $[1,2w]$ from upper $2w-S$ to lower $w-S$.
The inequality $w\ge S+1$ places both endpoints in their classes.
Raise its upper entries by $2T+2$. Prepend it to the old word with
every old entry raised by $w$.

The prefix lower block $[1,w]$ and old lower block $[w+1,w+T]$
are consecutive. The old upper block ends at $2T+w+2$ and the
prefix upper block is $[2T+w+3,2T+2w+2]$. The unique hole is
$T+V+w$. The prefix internal palette is
$[2T+3,2T+2w+1]$. Its final lower entry is $w-S$, while the old
first entry becomes $E+w$, so their join has color $2T+2$.
The old differences are unchanged. The first endpoint is
$2w-S+2T+2=2(T+w)-S+2$, and the last is $F+w$.
The missing color is unchanged since
$S+(T+w)-(U+w)+V=D$.
\end{proof}

\begin{lemma}[Asymmetric growth]\label{ok:asymmetric-growth}
From $K(S,U,T,V)$ and an integer
\[
 P\ge1,\qquad P\ge U+V-2S,
 \]
one obtains $K(S+P,U+2P,T+P,V)$.
\end{lemma}
\begin{proof}
Put
\[
 \ell=F+P-2T-2=P+2S-U-V+1.
 \]
The growth bound gives $\ell\ge1$. Since the old endpoint satisfies
$F\le2T+2$, we also have $\ell\le P$. Raise all old entries by $P$.
Append the even path on $[1,2P]$ from lower $\ell$ to upper
$\ell+P$, raising just its upper entries by $2T+2$.

The old labels lie in $[P+1,2T+P+2]$ except the translated hole.
The appended lower and upper labels fill respectively $[1,P]$ and
$[2T+P+3,2T+2P+2]$. They alternate across the new cut $T+P$.
The old colors are unchanged; the join has color
$(F+P)-\ell=2T+2$, and the new internal colors are
$[2T+3,2T+2P+1]$. The endpoints become $E+P,F+2P$.
These inventories have the stated parameters and the same missing
color $D$.
\end{proof}

\begin{lemma}[Reversal and class reflection]\label{ok:reversal}
\label{ok:reflection}
Reversing the word of a kernel gives the parameter transformation
\[
 J(S,U,T,V)=(U+V-2S-1,\ 2U+V-3S-1,\ T,\ V).
 \tag{J}
 \label{ok:J}
 \]
Complementing the labels separately within their classes, without
reversing the word, gives
\[
 R(S,U,T,V)=(T-S+1,\ 2T-U+2,\ T,\ T-V+3).
 \tag{R}
 \label{ok:R}
 \]
Both transformations are involutions on their realizable kernel
inventories.
\end{lemma}
\begin{proof}
For $J$, retain all labels. Substitution shows that the new first
endpoint is $F$, the new last endpoint is $E$, and the missing color
remains $D$. The classes and all other inventories are unchanged.
Applying the displayed affine transformation twice gives the identity.

For $R$ replace a lower entry $x$ by $T+1-x$ and an upper entry
$y$ by $3T+3-y$. This permutes each full class interval. The old
hole $T+V$ becomes $2T-V+3=T+(T-V+3)$. Every edge of color
$y-x$ becomes an edge of color $2T+2-(y-x)$. Its missing color is
therefore $2T+2-D$. The new endpoints are $3T+3-E$ and $3T+3-F$.
Substitution gives exactly \eqref{ok:R}; the label map is itself an
involution. Reversal is a separate operation: it is used in addition
to reflection only when explicitly stated.
\end{proof}

\subsubsection{Elementary seeds and a controlled-edge step}

\begin{lemma}[Three elementary sources]\label{ok:elementary-seeds}
The following kernels exist.
\begin{enumerate}
\item For $d\ge1$,
\[
 Z_d=(2d+2,1,2d+1,2,2d,\ldots,d,d+2)
       \quad\hbox{is }K(0,d,d,1).
 \]
\item For $r\ge2$, an even path on $[1,2r]$ from upper $r+2$
to lower $2$, followed by $2r+2$, is
$K(r,r,r,r+1)$.
\item For $h\ge2$, an even path on $[1,2h]$ from lower $2$ to
upper $h+2$, preceded by $2h+2$, is $K(0,0,h,h+1)$.
\end{enumerate}
\end{lemma}
\begin{proof}
In the first word the upper entries are $2d+2,2d+1,\ldots,d+2$;
its consecutive differences are $2d+1,2d,\ldots,2$.
Thus its label hole is $d+1$, color hole is $1$, and its endpoints
are $2d+2,d+2$, proving the first assertion.

In the second assertion the even path supplies $[1,2r-1]$ and
the appended edge has color $2r$. The omitted label is $2r+1$,
the missing color is $2r+1$, and the endpoints are $r+2,2r+2$.
In the third assertion the prepended edge has color $2h$.
The omitted label and missing color are both $2h+1$, and the
endpoints are $2h+2,h+2$. These agree with the kernel formulas.
The even-path requests have $2\le r$ and $2\le h$ respectively,
so all their prescribed endpoints are valid.
\end{proof}

\begin{lemma}[A controlled-edge enlargement]\label{ok:controlled-step}
Suppose a kernel $K(S,U,T,V)$ contains the edge $(2,T+1)$ and
$T\ge2$. There is a kernel
\[
 K(S+2,U+4,T+4,V+2)
 \]
which contains the new controlled edge $(2,T+5)$.
\end{lemma}
\begin{proof}
Raise old lower labels by $2$ and old upper labels by $6$.
The mapped cut is $T+4$. The old label blocks are $[3,T+2]$ and
$[T+7,2T+8]$ with the old hole raised by $6$. The fresh labels are
\[
 1,2,T+3,T+4;\qquad T+5,T+6,2T+9,2T+10.
 \]
The old controlled edge maps to $(4,T+7)$, with color $T+3$.
Replace this edge, using the reverse order if necessary, by
\begin{equation}
 (4,2T+10,1,2T+9,2,T+5,T+4,T+6,T+3,T+7).
 \label{ok:controlled-chain}
\end{equation}
Its nine colors in order are
\[
 (2T+6,2T+9,2T+8,2T+7,T+3,1,2,3,4).
 \]
The old controlled color is retained. Every other old color rose by
$4$, so the new colors added are exactly
$[1,4]\cup[2T+6,2T+9]$. This gives the new palette with missing
color $D+4$. All fresh entries occur once and cross the new cut.
The hole and both endpoints have risen by $6$.
These are precisely the changes specified by the new parameters.
The replacement contains $(2,T+5)$. Replacing one edge of a path
by a path through eight fresh vertices preserves a single path, so
there is no separate connectivity assumption.
\end{proof}

\subsection{All positive odd three-arm spiders}
\label{construction:all-odd-alpha3}

We first give the constructions for minimum arm at least five. The
two smaller minima and the single negative profile are treated at the
end of the section.

\subsubsection{Equal and repeated largest arms}

\begin{proposition}[Three equal odd arms]\label{ok:equal-odd-three}
For every $r\ge2$, $\Sp{(2r+1)^3}$ is root-zero $\alpha$, with
index $3r$, and its first arm is the maximum Walecki word.
\end{proposition}
\begin{proof}
At $r=2$ use
\[
 (15,1,14,2,13)\mid(10,3,8,4,12)\mid(9,6,7,5,11);
 \]
at $r=3$ use
\[
 (21,1,20,2,19,3,18)\mid
 (14,5,15,7,12,6,17)\mid(13,9,10,8,11,4,16).
 \]
Their lower labels are respectively $[1,6]$ and $[1,9]$,
and direct successive subtraction gives their full color intervals.

Now let $r\ge4$, $N=6r+3$, and $K=3r$.
For upper list $(U_0,\ldots,U_r)$ and lower list
$(L_0,\ldots,L_{r-1})$, interleaving means the word
$(U_0,L_0,U_1,\ldots,L_{r-1},U_r)$.
Take the first arm to be
\[
 W=(N,1,N-1,2,\ldots,r,N-r).
 \]
Its lower block is $[1,r]$, upper block $[5r+3,6r+3]$, and
palette $[4r+3,6r+3]$. For arm $A$ interleave
\begin{align*}
 U_A&=(4r+1,4r-1,4r+3,4r+4,\ldots,5r+1),\\
 L_A&=(2r-1,2r,2r-2,2r-3,\ldots,r+1).
\end{align*}
Its root color is $4r+1$, its first four internal colors are
$2r+2,2r,2r-1,2r+3$, and the remaining colors are consecutively
$2r+5,\ldots,4r$. Hence
\begin{equation}
 C_A=\{4r+1,2r-1,2r,2r+2,2r+3\}\cup[2r+5,4r].
 \label{ok:equal-three-A}
\end{equation}

If $r$ is even, take
\[
 L_C=(2r+1,\ldots,3r-3,3r-1,3r,3r-2).
 \]
The upper list begins with $4r+2$, then appends
$(4r-3-2j,4r-2j)$ for $0\le j\le r/2-3$, and finally appends
$(3r+2,3r+1,3r+4,5r+2)$. An empty loop is omitted.
The root color and first internal color are $4r+2,2r+1$.
Each loop packet contributes
\[
 2r-4-4j,\quad2r-5-4j,\quad2r-2-4j,\quad2r-3-4j.
 \]
Their union is $[7,2r-2]$. The seven final internal colors are
$5,3,2,1,4,6,2r+4$. Thus the complete palette is
\begin{equation}
 C_C=[1,2r-2]\cup\{2r+1,2r+4,4r+2\}.
 \label{ok:equal-three-C}
\end{equation}
If $r$ is odd, so $r\ge5$, use instead
\[
 L_C=(2r+1,\ldots,3r-3,3r,3r-1,3r-2).
 \]
The upper list begins with $4r+2$, then appends the same pairs
for $0\le j\le(r-7)/2$, then appends
$(3r+3,3r+5,3r+1,3r+2,5r+2)$.
The packets give $[9,2r-2]$ and the final nine colors are
$7,6,8,5,1,2,3,4,2r+4$, giving the same
\eqref{ok:equal-three-C}.

The lower lists partition $[1,3r]$. Arm $A$ has upper set
$\{4r-1,4r+1\}\cup[4r+3,5r+1]$; arm $C$ has upper set
\[
 ([3r+1,4r+2]\setminus\{4r-1,4r+1\})\cup\{5r+2\}.
 \]
Together with the upper block of $W$ these partition
$[3r+1,6r+3]$. Every edge crosses $3r$.
The three disjoint palettes
\eqref{ok:equal-three-A}, \eqref{ok:equal-three-C}, and that of
$W$ partition $[1,6r+3]$. This proves every inventory.
\end{proof}

\begin{lemma}[Repeated-arm kernel sources]\label{ok:repeated-kernels}
For $s\ge2$ and $t>s$, there is a kernel $K(s,t,t,v)$ with
$1\le v\le t$. More precisely:
\begin{enumerate}
\item If $s<t\le2s-2$, one can take $v=s+1$.
\item If $t\ge2s+1$, one can take $v=s+1$.
\item If $t=2s-1$, one can take $v=s$.
\item If $t=2s$ and $s\ge3$, one can take $v=s+3$.
At $(s,t)=(2,4)$ one can take $v=2$.
\end{enumerate}
\end{lemma}
\begin{proof}
In the first case put $d=t-s$, so $1\le d\le s-2$.
Apply diagonal growth with $P=s$ to $Z_d=K(0,d,d,1)$.
The required inequality is $s\ge d+2$, and the result is
$K(s,t,t,s+1)$.

In the second case start with $K(s,s,s,s+1)$ from
Lemma~\ref{ok:elementary-seeds} and apply carrier prefix with
$w=t-s\ge s+1$. This gives the same standard parameters.
The apparent failure $s+1>s$ in the initial source is harmless:
completion is used only after the prefix, when $s+1\le t$.

For $t=2s-1$ the desired label hole is $3s-1$, missing color
$2s$, and endpoints $3s,3s+2$. Two initial words are
\begin{align*}
 P_2&=(6,3,4,2,7,1,8),\\
 P_3&=(9,4,7,5,6,2,10,3,12,1,11).
\end{align*}
Both contain $(2,t+1)$ and have exactly the stated inventories.
Lemma~\ref{ok:controlled-step} changes $(s,t,t,s)$ to
$(s+2,t+4,t+4,s+2)$. Since $t+4=2(s+2)-1$, its two initial
words give every $s\ge2$.

For $t=2s$, $s\ge3$, take the standard source
$K(s-1,2s,2s,s)$ from the second case. First reflect within
classes and then reverse its word. The composite $J\circ R$ gives
$K(s,2s,2s,s+3)$: its hole is $3s+3$, missing color $2s+3$,
and ordered endpoints $3s+2,3s$. Here $s+3\le2s$.
For the isolated case $(2,4)$ the word
\[
 (8,3,5,4,7,1,10,2,9)
 \]
is $K(2,4,4,2)$; its hole is $6$, missing color $4$, and
endpoints $8,9$. These four cases exhaust $t>s$.
\end{proof}

\begin{corollary}\label{ok:repeated-alpha}
For odd $5\le a\le b$, $\Sp{a,b,b}$ is root-zero $\alpha$ and
admits a maximum Walecki $a$-arm. For $a<b$ this follows from
Lemma~\ref{ok:completion} and Lemma~\ref{ok:repeated-kernels};
for $a=b$ it follows from Proposition~\ref{ok:equal-odd-three}.
\end{corollary}

\subsubsection{The interior central cone and its arithmetic progressions}

\begin{lemma}[The non-AP central kernel]\label{ok:central-kernel}
Let
\[
 1\le s\le t\le u\le2s-2,\qquad u>s,\qquad u\ne2t-s.
 \]
There is a kernel $K(s,u,t,s+t-u+1)$.
\end{lemma}
\begin{proof}
Put
\[
 S=u-t,\quad T=u-s,\quad U=2u-s-t=S+T,\quad P=s+t-u.
 \]
Then $0\le S\le T$, $T\ge1$, $U\ge1$, and
$P\ge T-S+2$. Apply Lemma~\ref{common:odd-path} to the
$2T+1$-vertex path with upper first entry $2T-S+1$.
Its other endpoint is $T+S+1$. The first entry is between
$T+1$ and $2T+1$. The excluded request is
\[
 2(2T-S)=3T\quad\Longleftrightarrow\quad T=2S
 \quad\Longleftrightarrow\quad u=2t-s,
 \]
which is excluded in the hypothesis. The normalization in
Lemma~\ref{common:odd-path} also handles $S=T$, when the prescribed
entry is the first upper entry.

Raise every upper entry of this ordinary path by one. The new label
hole is $T+1$, the missing color is $1$, and its endpoints are
$2T-S+2,T+S+2$. Thus it is $K(S,U,T,1)$.
Diagonal growth by $P$ gives $K(s,u,t,P+1)$.
Finally $1\le P+1\le u$, since $U\ge1$ and $u=U+P$.
\end{proof}

\begin{lemma}[An explicit near-maximum-hole word]\label{ok:near-max-word}
For $q\ge2$ the following concatenation is a path word:
\[
 (3q+1),\quad
 (q+j,3q-j)_{j=0}^{q-1},\quad
 (2q,4q+2),\quad
 (j,4q+1-j)_{j=1}^{q-1}.
 \tag{*}\label{ok:near-max-formula}
 \]
Its lower labels are $[1,2q]$, its upper labels are
$[2q+1,4q+2]\setminus\{4q+1\}$, its colors are
$[1,4q+1]\setminus\{4q\}$, and its endpoints are $3q+1,3q+2$.
\end{lemma}
\begin{proof}
The lower packets are $[q,2q-1]$, $\{2q\}$, and $[1,q-1]$.
The upper packets are $\{3q+1\}$, $[2q+1,3q]$,
$\{4q+2\}$, and $[3q+2,4q]$. These are disjoint and give
the stated label sets.

Through the entry $2q$, the differences are exactly $[1,2q+1]$:
the first is $2q+1$, the first loop supplies alternating descending
even and odd differences down to $2$, and the following edge gives
$1$. The next edge gives $2q+2$. The remaining edges give all of
$[2q+3,4q+1]$ except $4q$: the within-pair differences are
$4q+1-2j$, and the joins give the intervening even differences.
All edges cross $2q$. The first and last entries are as asserted.
\end{proof}

\begin{lemma}[Pinned central AP kernels]\label{ok:ap-central}
If $d\ge2$ and $s\ge2d+2$, there is
\[
 K(s,s+2d,s+d,s-d+2).
 \]
For $d=1$ and every $s\ge2$ there is also a kernel completing
$\Sp{2s+1,2s+3,2s+5}$.
\end{lemma}
\begin{proof}
In \eqref{ok:near-max-formula} put $q=d$. Reflect its lower
labels by $2d+1-x$ and upper labels by $6d+3-x$. Its new hole
is $2d+2$, missing color $2$, and endpoints $3d+2,3d+1$.
These are exactly $K(d,3d,2d,2)$. Equivalently the reflected word
is the following explicit concatenation:
\[
 (3d+2),\quad
 (d+1-j,3d+3+j)_{j=0}^{d-1},\quad
 (1,2d+1),\quad
 (2d+1-j,2d+2+j)_{j=1}^{d-1}.
 \]
Apply diagonal growth with $P=s-d\ge d+2$. The result has
parameters $(s,s+2d,s+d,s-d+2)$, and its final hole parameter
is between $1$ and $s+2d$.

For $d=1$ the following three literal kernels supply $s=2,3,4$:
\begin{align*}
 K(2,4,3,2)&=(6,1,8,2,4,3,7),\\
 K(3,5,4,2)&=(7,4,5,3,8,2,9,1,10),\\
 K(4,6,5,4)&=(8,2,12,1,10,5,6,4,7,3,11).
\end{align*}
Their respective label holes are $5,6,9$ and missing colors
$3,4,7$, so direct subtraction checks all three.
For $s\ge5$, diagonally grow the first word by $P=s-2\ge3$.
It gives $K(s,s+2,s+1,s)$, with $1\le s\le s+2$.
The completion lemma applies in every case, with the Walecki
short arm retained.
\end{proof}

\begin{corollary}\label{ok:interior-cone}
If $5\le a\le b\le c\le2a-5$ are odd, then
$\Sp{a,b,c}$ is root-zero $\alpha$.
\end{corollary}
\begin{proof}
Write $a=2s+1$, $b=2t+1$, $c=2u+1$.
The equal case uses Proposition~\ref{ok:equal-odd-three}.
The non-AP case uses Lemmas~\ref{ok:central-kernel}
and~\ref{ok:completion}. In the remaining case put
$t=s+d$, $u=s+2d$, $d\ge1$.
The cone condition implies $s\ge2d+2$, so
Lemma~\ref{ok:ap-central} applies. This covers the complete
arithmetic-progression exception of the ordinary odd-path request.
\end{proof}

\subsubsection{The two central boundary layers}

\begin{lemma}[Odd-prefix boundary kernels]\label{ok:odd-prefix-boundaries}
For $\epsilon\in\{0,1\}$, suppose
\[
 u=2s-1+\epsilon,\qquad
 s\ge3+\epsilon,\qquad s\le t\le2s-3-\epsilon.
 \]
Then there is $K(s,u,t,t-s+3)$.
\end{lemma}
\begin{proof}
Put
\[
 T=s-1,\quad P=t-s+1,\quad L=2s-t-2-\epsilon,\quad v=P+2.
 \]
Here $T,P\ge1$ and $1\le L\le T$.
Take the even path on $[1,2T]$ from lower $L$ to upper $T+L$,
raising its upper entries by $2P+2$. Prepend the odd Walecki
word on $[1,2P+1]$, starting at $2P+1$ and ending at $P+1$,
translated by $T$.

The lower blocks are $[1,T]$ and $[T+1,t]$. The upper blocks
are $[t+1,T+2P+1]$ and $[T+2P+3,2t+2]$.
The label hole is $T+2P+2=t+v$.
The prefix colors are $[1,2P]$; the old even path contributes
$[2P+3,2t+1]$; and their join has color
\[
 (T+P+1)-L=2P+1+\epsilon.
 \]
Thus the sole missing color is $2P+2-\epsilon=s+t-u+v$.
The endpoints are
\[
 T+2P+1=2t-s+2,\qquad
 T+L+2P+2=t+s+1-\epsilon.
 \]
The second is $2t+2s-u-v+3$, proving the claimed kernel.
The inequalities also give $1\le v\le u$ when completion is wanted.
There is no odd-order prescribed endpoint in this argument: the odd
prefix is an explicit word.
\end{proof}

\begin{lemma}[The gap-two boundary kernels]\label{ok:gap-two-kernels}
For every $s\ge4$ there is
\[
 K(s,2s-1,2s-2,s+3).
 \]
Consequently $\Sp{a,2a-5,2a-3}$ is root-zero $\alpha$ for
every odd $a\ge5$.
\end{lemma}
\begin{proof}
For the displayed kernel the carrier is $T=2s-2$, the hole is
$3s+1$, missing color $2s+2$, and endpoints $3s-2,3s-3$.
At $s=4$ and $s=7$ use the controlled initial words
\begin{align*}
 P_4={}&(10,4,11,3,12,1,14,2,7,6,8,5,9),\\
 P_7={}&(19,7,20,6,21,4,25,1,26,3,23,5,24,2,\\
       &\hspace{9mm}13,12,14,11,15,10,16,9,17,8,18).
\end{align*}
They contain $(2,T+1)$ and have exactly the inventories just stated.
Lemma~\ref{ok:controlled-step} supplies $s\mapsto s+2$.
The $s=4$ word covers all even $s\ge4$, and the $s=7$ word
all odd $s\ge7$. The one remaining value, $s=5$, is supplied by
\[
 (13,6,9,8,10,4,18,1,17,2,15,5,14,3,11,7,12).
 \]
No controlled edge is required in this isolated word.
Since $s+3\le2s-1$ for $s\ge4$, completion is valid.

At $a=7$ the profile is $(7,9,11)$, supplied by the
$s=3,d=1$ word in Lemma~\ref{ok:ap-central}.
At $a=5$ it is $(5,5,7)$ and uses the complete
minimum-five theorem in Section~\ref{app:small-alpha}.
\end{proof}

\begin{lemma}[The reverse boundary]\label{ok:reverse-boundary}
For $s\ge2$ and $s\le t\le2s$, the spider
$\Sp{2s+1,2t+1,4s+1}$ is root-zero $\alpha$.
For $t<2s$ it has a completing kernel $K(s,2s,t,t-s+1)$.
\end{lemma}
\begin{proof}
First suppose $t<2s$ and put $r=t-s$, so $0\le r\le s-1$.
We construct the standard source $K(r,t,t,r+1)$.
For $r=0$ use $Z_t$. For $r=1$ use
\[
 B_3=(7,3,4,2,8,1,6),\qquad
 B_4=(9,3,5,4,8,1,10,2,7)
 \tag{B}\label{ok:B-small}
 \]
when $t=3,4$, respectively. They are $K(1,t,t,2)$.
For $t\ge5$, apply carrier prefix of size $w=t-3\ge2$ to
$B_3$; its even upper endpoint is $2w-1$.
For $r\ge2$, the inequality $t=s+r\ge2r+1$ permits the
standard short-wedge source in Lemma~\ref{ok:repeated-kernels}.
Reversal $J$ gives exactly $K(s,2s,t,r+1)$.
Its hole parameter satisfies $1\le r+1\le s\le2s$.
When $t=2s$, use the twice-carrier case of
Lemma~\ref{ok:repeated-kernels}, including its literal $(2,4)$
exception. Completion proves the assertion.
\end{proof}

\begin{proposition}[The complete central cone]\label{ok:central-cone}
For odd $5\le a\le b\le c\le2a-1$, the spider
$\Sp{a,b,c}$ is root-zero $\alpha$.
For $a\ge7$ these constructions retain the maximum Walecki $a$-arm.
\end{proposition}
\begin{proof}
The interior $c\le2a-5$ is Corollary~\ref{ok:interior-cone}.
At $c=2a-3$, either $b=c$, handled by the repeated-arm theorem,
or $b=c-2$, handled by Lemma~\ref{ok:gap-two-kernels}, or
$b\le c-4$, handled by $\epsilon=0$ in
Lemma~\ref{ok:odd-prefix-boundaries}.
At $c=2a-1$ use Lemma~\ref{ok:reverse-boundary}.
All three possibilities at each layer exhaust odd $b$.
For the small minimum $a=5$ the complete minimum-five theorem also
supplies any branch outside those displayed parameter domains.
Every construction for $a\ge7$ is a completion or the explicitly
pinned equal-arm construction. No Walecki pin is asserted for the
separate minimum-five input.
\end{proof}

\subsubsection{The far cone, including its resonance}

\begin{lemma}[Even extension at an upper leaf]\label{ok:upper-leaf-extension}
Suppose a root-zero $\alpha$-labeled tree has index $K$ and an
upper leaf labeled $K+e$. Any $2z$ vertices, with $z\ge e$, can be
appended to that leaf while preserving root-zero $\alpha$.
The new index is $K+z$ and the new leaf has label $K+z+e$.
\end{lemma}
\begin{proof}
Raise every old upper label by $2z$, leaving old lower labels fixed.
On the new middle interval use an even path on $[1,2z]$
from lower $e$ to upper $z+e$, adding $K$ to its entries.
Join its first entry to the former leaf. The join has color
$(K+e+2z)-(K+e)=2z$, the new internal colors are
$[1,2z-1]$, and all old colors have risen by $2z$.
The middle labels fill $[K+1,K+2z]$ exactly. Every edge crosses
$K+z$, and the asserted new leaf and distance follow.
\end{proof}

\begin{proposition}[The full odd far cone]\label{ok:far-cone}
If $7\le a\le b\le c$ are odd and $c\ge a+b-1$, then
$\Sp{a,b,c}$ is root-zero $\alpha$.
\end{proposition}
\begin{proof}
Write $a=2s+1$, $b=2t+1$, with $s\ge3$.
The rooted pair on arms $a,b$ is root-zero $\alpha$ of index
$s+t$, by Corollary~\ref{common:alpha-pair}.
Since $c>2(s+t)$, Lemma~\ref{common:alpha-extension}
attaches the $c$-arm except at
\[
 c_*=4(s+t)+1=2a+2b-3.
 \]
We now construct that value, rather than assume attachment at its
excluded endpoint.

Take the pinned repeated-arm construction $\Sp{a,b,b}$ from
Corollary~\ref{ok:repeated-alpha}, including its equal case.
Put $K=s+2t$ and $N=2K+3=2s+4t+3$.
All three leaves are upper. If their labels sum to $L$, summing the
positive edge differences by degrees gives
\[
 \frac{N(N+1)}2
   =2\sum_{j=K+1}^N j-2\sum_{j=1}^Kj-L,
 \]
because nonleaf nonroot vertices have degree two and the hub is zero.
Consequently
\begin{equation}
 L=\frac{N(N+1)}2-2K(K+1)=5K+6.
 \label{ok:signed-leaf-sum}
\end{equation}
The Walecki short leaf is $N-s=s+4t+3$.
The two long leaf labels therefore sum to $4s+6t+3$.
Their smaller label $U$ satisfies
\[
 U\le2s+3t+1,\qquad e=U-K\le s+t+1.
 \]
Choose that long arm and append $2(2s+t)$ vertices.
Since $2s+t\ge s+t+1\ge e$, Lemma~\ref{ok:upper-leaf-extension}
applies. Its length becomes
\[
 b+2(2s+t)=4s+4t+1=c_*.
 \]
This supplies the resonance.
The extension raises all old upper labels uniformly, so the short
arm remains the maximum Walecki word. The unextended lower labels
and all new middle labels stay below its final upper label.
\end{proof}

\subsubsection{The complete short wedge}

\begin{lemma}[Unequal short-wedge kernels]\label{ok:unequal-wedge-kernel}
Let $s\ge2$, $u\ge2s+2$, and $0\le\delta\le s-1$; put
$t=u+\delta$. There is a completing kernel $K(s,u,t,v)$, with
\[
 v=\begin{cases}
 s+\delta+3,&\delta\ne s-3,\\
 2s-1,&\delta=s-3.
 \end{cases}
 \]
\end{lemma}
\begin{proof}
In the nonexceptional case put $r=s+\delta+1$.
Apply the odd-path interface on $[1,2r+1]$ at upper endpoint
$2r-s+2$; its other endpoint is $2s+\delta+1$.
Both entries lie in the upper class. The sole excluded request is
\[
 2(2r-s+1)=3r
 \quad\Longleftrightarrow\quad r=2s-2
 \quad\Longleftrightarrow\quad\delta=s-3.
 \]
Thus the request is permitted here.
The word itself, with a new unused possible maximum $2r+2$,
is $K(s,s+1,r,s+\delta+3)$:
its hole is $r+v=2r+2$ and its missing color is $2r+1$.

If $\delta=s-3$, then $s\ge3$. Put $q=s-1$ in
Lemma~\ref{ok:near-max-word}. Its word is
$K(s,s+1,2s-2,2s-1)$, with missing color $4s-4$.
This is the needed replacement for the forbidden ordinary path.
In both cases apply carrier prefix with
$w=u-s-1\ge s+1$. The carrier becomes $t=u+\delta$,
the other half becomes $u$, and $v$ is unchanged.
Since $v\le2s+2\le u$ in the first case and $v=2s-1<u$
in the second, completion is allowed.
\end{proof}

\begin{proposition}[Both adjacent short-wedge boundaries]
\label{ok:short-wedges}
For odd $5\le a\le b\le c$ with $b\ge2a-1$,
$\Sp{a,b,c}$ is root-zero $\alpha$.
\end{proposition}
\begin{proof}
We can dispose of $a=5$ using Section~\ref{app:small-alpha};
write $a=2s+1$ with $s\ge3$ for the far-cone appeals below.

First let $b\ge2a+3$, so its half $u\ge2s+2$.
Write the largest half as $t=u+\delta$.
If $0\le\delta\le s-1$, use
Lemma~\ref{ok:unequal-wedge-kernel} and completion.
If $\delta\ge s$, then $c\ge a+b-1$ and
Proposition~\ref{ok:far-cone} applies.

Next let $b=2a+1$, with half $u=2s+1$.
For $2s+1\le t\le3s$ take the $\epsilon=0$ odd-prefix kernel
whose source parameters are
\[
 S_0=t-s+1,\quad U_0=2S_0-1,\quad T_0=t,\quad V_0=s+2.
 \]
Indeed $S_0\ge3$, $S_0\le t\le2S_0-3$,
as required by Lemma~\ref{ok:odd-prefix-boundaries}.
Reflecting it gives
\[
 K(s,2s+1,t,t-s+1).
 \]
The final hole parameter is at most $2s+1$ exactly because
$t\le3s$. For $t\ge3s+1$ the far cone applies.

Finally let $b=2a-1$, with half $u=2s$.
If $t=2s$, use the repeated-arm theorem.
If $t=2s+1$, reflect the gap-two source
$K(s+1,2s+1,2s,s+4)$ of
Lemma~\ref{ok:gap-two-kernels}. The result is
$K(s,2s+1,2s,s-1)$.
Completing it gives the desired unordered arm triple, with its two
long arms interchanged. The condition $s\ge3$ ensures $s-1\ge1$.
For $2s+2\le t\le3s-1$, use the $\epsilon=1$ odd-prefix source
with $S_0=t-s+1$, $U_0=2S_0$, $T_0=t$, $V_0=s+2$.
Its source inequalities hold, and reflection gives
$K(s,2s,t,t-s+1)$. Its final hole condition is exactly
$t\le3s-1$. For $t\ge3s$, the far cone applies.
All equality, gap-two, near, and far cases have now been included.
\end{proof}

\subsubsection{The last transition and its pinned AP ray}

\begin{lemma}[Transition kernels]\label{ok:transition-kernels}
The following kernels exist:
\begin{align}
 K(s,2s+g,s+g+h,h+1),
 &\quad g\ge1,\ h\ge2,\ s\ge g+h+1; \label{ok:transition-high}\\
 K(s,2s+g,s+g+1,3),
 &\quad g\ge1,\ s\ge g+3. \label{ok:transition-one}
\end{align}
\end{lemma}
\begin{proof}
For \eqref{ok:transition-high}, use the elementary pinned source
$K(0,0,h,h+1)$ from Lemma~\ref{ok:elementary-seeds}.
Carrier prefix of size $g$ gives $K(0,g,g+h,h+1)$.
The inequality $s\ge g+h+1=U+V$ permits asymmetric growth
by $P=s$, giving the asserted kernel. Its ordinary even appended
path has lower endpoint $s-g-h\ge1$ and upper endpoint
$2s-g-h\le2s$.

For \eqref{ok:transition-one}, the literal word
\[
 (6,1,4,2,3)=K(0,1,2,3)
 \]
has label hole $5$, missing color $4$, and endpoints $6,3$.
If $g>1$, use carrier prefix of size $g-1$; if $g=1$ do
nothing. The resulting kernel is $K(0,g,g+1,3)$.
Asymmetric growth by $s\ge g+3=U+V$ proves the second assertion.
The final hole parameters $h+1$ and $3$ lie in $[1,2s+g]$,
so either word can be completed.
\end{proof}

\begin{lemma}[The pinned limiting AP ray]\label{ok:pinned-ap-ray}
For every $s\ge3$ there is $K(s,3s-2,2s-1,2)$.
Its completion is a root-zero $\alpha$-labeling of
$\Sp{2s+1,4s-1,6s-3}$ with a maximum Walecki short arm.
\end{lemma}
\begin{proof}
The required word has label interval $[1,4s]\setminus\{2s+1\}$,
cut $2s-1$, palette $[1,4s-1]\setminus\{3\}$, and endpoints
$3s,3s+1$. At $s=3$ use
\[
 (9,5,6,4,11,3,8,2,12,1,10).
 \]
At $s=4,5$ use the two controlled initial words
\begin{align*}
 Q_4={}&(12,6,8,7,11,2,16,1,14,4,15,3,10,5,13),\\
 Q_5={}&(15,6,12,8,10,9,14,7,17,5,13,2,20,1,18,4,19,3,16).
\end{align*}
They contain the controlled edge $(2,2s+3)$.
These finite assertions follow by listing their absent label, endpoints,
and successive absolute differences: the latter give every integer
from $1$ to $4s-1$ except $3$.

For $s\ge4$ with a controlled word, raise lower entries by $2$
and upper entries by $6$. The mapped blocks are $[3,2s+1]$ and
$[2s+6,4s+6]\setminus\{2s+7\}$; the palette is
$[5,4s+3]\setminus\{7\}$. The controlled edge becomes
$(4,2s+9)$, of color $2s+5$. Replace it by
\[
 (4,4s+8,1,4s+7,2,2s+7,2s+3,2s+4,2s+2,2s+9),
 \]
or its reversal as dictated by the old edge orientation.
Its eight interior entries are the fresh lower labels
$1,2,2s+2,2s+3$ and upper labels
$2s+4,2s+7,4s+7,4s+8$.
Its colors in order are
\[
 (4s+4,4s+7,4s+6,4s+5,2s+5,4,1,2,7).
 \]
Thus the controlled color is retained; the added colors are
$1,2,4,7$ and $[4s+4,4s+7]$.
Together with the mapped old palette these give
$[1,4s+7]\setminus\{3\}$. The target hole is $2s+5$,
the new endpoints are $3s+6,3s+7$, and the new controlled edge
is $(2,2s+7)$. All edges cross $2s+3$.
The replacement uses fresh entries once and preserves the path topology.
This is the induction $s\mapsto s+2$; the two bases cover both
parities. Finally $1\le2\le3s-2$ permits completion.
\end{proof}

\begin{theorem}[Complete positive odd classification]\label{ok:all-odd-theorem}
A positive three-arm spider with all arm lengths odd is root-zero
$\alpha$ if and only if its profile is not $(3,3,3)$.
\end{theorem}
\begin{proof}
Order its lengths $a\le b\le c$.
First let $a\ge7$ and write $a=2s+1$, $b=2t+1$, $c=2u+1$.
The complete central cone $u\le2s$, complete short wedge
$t\ge2s$, and far cone $u\ge s+t$ have already been proved.
Outside their union,
\[
 u\ge2s+1,\qquad t\le2s-1,\qquad u\le s+t-1.
 \]
Put $g=u-2s$ and $h=s+t-u$. Then $g,h\ge1$ and
$t=s+g+h\le2s-1$, or $s\ge g+h+1$.
For $h\ge2$ use \eqref{ok:transition-high}.
For $h=1$ and $s\ge g+3$ use \eqref{ok:transition-one}.
The only remaining case is $h=1$, $s=g+2$, giving
$(a,b,c)=(2s+1,4s-1,6s-3)$; it is
Lemma~\ref{ok:pinned-ap-ray}. These alternatives exhaust every
positive size, without an upper bound or a total-size congruence.
For $a=5$, the equal profile $(5,5,5)$ is
Proposition~\ref{ok:equal-odd-three}; every other profile has
$c\ge7$ and is supplied by the minimum-five theorem in
Section~\ref{app:small-alpha}.

If $a=1$, the rooted odd pair $\Sp{b,c}$ is root-zero
$\alpha$ by Corollary~\ref{common:alpha-pair}; its only exceptional
rooted pair $(2,2)$ cannot occur. Adding a maximum-labeled unit
leaf preserves its index and proves the assertion.
If $a=3$ and $c\ge7$, use the complete minimum-three theorem
in Section~\ref{app:small-alpha}. If $c<7$, the two positive
profiles have the following explicit words:
\begin{align*}
 \Sp{3,3,5}:&
 \quad(11,1,10)\mid(6,2,9)\mid(8,3,5,4,7),\\
 \Sp{3,5,5}:&
 \quad(13,1,12)\mid(10,4,6,5,8)\mid(9,2,7,3,11).
\end{align*}
They have cuts $4,5$, respectively, consecutive labels, and
consecutive color inventories. The sole remaining profile is
$(3,3,3)$.

For this negative case a root-zero $\alpha$-labeling would have
cut $3$. Its depth-two labels would be $1,2,3$ and its first and
last labels $4,\ldots,9$. Equation~\eqref{ok:signed-leaf-sum}
at $K=3,N=9$ forces the leaf labels to sum to $21$.
Label $9$ must be a first label, since only an edge incident with
zero can have color $9$. The only three-element subset of
$[4,8]$ with sum $21$ is $\{6,7,8\}$; hence the first labels
are $4,5,9$.
Color $8$ forces $9$--$1$. Color $7$ then forces $1$--$8$:
the alternative $9$--$2$ would give degree three to the vertex
labeled $9$, while $7,8$ cannot be first labels.
Thus one entire arm is $(9,1,8)$.
The arm starting with $4$ cannot use middle $2$, because its final
edge to $6$ or $7$ would repeat root color $4$ or $5$.
It must therefore be $(4,3,6)$, since leaf $7$ would again repeat
color $4$. The last arm is then $(5,2,7)$ and repeats colors
$3,5$. This contradiction proves the exception.
The conclusion concerns root-zero $\alpha$ labeling; it does not
assert that $\Sp{3,3,3}$ is nongraceful.
\end{proof}

\subsection{Odd-pair offsets extracted from kernels}
\label{construction:odd-pair-kernels}

The basic pair offsets, including the low, high, reverse, equal, and
maximal constructions, are proved in Section~\ref{app:pair-offsets}.
The purpose of this section is to prove the additional odd-pair
intervals used in the eight-arm argument. Each statement is sufficient;
their union is not asserted to classify every possible offset.

\begin{lemma}[Kernel-to-pair extraction]\label{ok:pair-extraction}
Let $p$ be $K(r,u,t,V)$, with $t,u\ge1$ and $1\le V\le u$.
Then
\[
 k=t+u+V-1\ \in\ \Om(\Sp{2t+1,2u+1}).
 \]
No ordering assumption between $t$ and $u$ is needed.
\end{lemma}
\begin{proof}
Put $m=2t+2u+2$, $H=t+3-V$, and $\ell=u-V+1$.
Complement the kernel entries in $2t+3$. Its lower set is
$[1,t+2]\setminus\{H\}$, its upper set $[t+3,2t+2]$,
and its endpoints are $r+1,u-2r+V$.
Map its lower entries by $x\mapsto x+k$ and upper entries by
$x\mapsto x+k+2u$ to form arm $A$.
The even path on $[1,2u]$ from lower $\ell$ to upper
$u+\ell$ exists since $1\le\ell\le u$.
Translate it by $k+t+2$ and append $k+H$ to form arm $B$.

The two words have lengths $2t+1,2u+1$. Their nonroot labels
partition the following blocks, the indicated hole being filled last:
\begin{align*}
 &[k+1,k+t+2]\setminus\{k+H\},\\
 &[k+t+3,k+t+2u+2],\\
 &[k+t+2u+3,k+m],\qquad\{k+H\}.
\end{align*}
Thus their union is exactly $[k+1,k+m]$.
Complementation preserves the old differences, so the internal
colors of $A$ are
$[2u+1,m-1]\setminus\{D+2u\}$.
Its root color is $k+r+1=D+2u$.
The old even path in $B$ supplies $[1,2u-1]$,
its appended edge has color
\[
 (k+t+2+u+\ell)-(k+H)=2u,
 \]
and its root color is $k+t+2+\ell=m$.
Hence every color in $[1,m]$ appears once.
\end{proof}

\subsubsection{An extended interval for two equal odd arms}

\begin{lemma}[A variable-hole repeated kernel]\label{ok:variable-hole}
For $s\ge2$ and $1\le V\le s$, some integer $r\ge0$ admits
$K(r,s,s,V)$.
\end{lemma}
\begin{proof}
Normally put $r=V-1$, so the desired hole parameter is $r+1$.
For $r=0$, use $Z_s$. For $r=1$, the words in
\eqref{ok:B-small} apply at $s=3,4$, and carrier prefix of
size $s-3\ge2$ from $B_3$ applies at $s\ge5$.
For $r\ge2$, diagonal growth from $Z_{s-r}$ gives the standard
kernel if $s\le2r-2$, and carrier prefix from
$K(r,r,r,r+1)$ gives it if $s\ge2r+1$.
Here $r=V-1\le s-1$, so the diagonal source has positive gap
whenever it is used.

The two omitted carriers are $s=2(V-1)-1$ and $s=2(V-1)$.
In these cases change to $r=V-2$.
If $r=0$, necessarily $(s,V)=(2,2)$, and use
\[
 (6,1,5,2,3)=K(0,2,2,2).
 \]
If $r\ge1$, put $S=s-r+1$.
The two carrier alternatives give $S\ge3$ and
$S\le s\le2S-3$.
The $\epsilon=0$ odd-prefix source
of Lemma~\ref{ok:odd-prefix-boundaries} has parameters
$K(S,2S-1,s,r+2)$.
Reversing it by $J$ gives $K(r,s,s,r+2)$,
which has the desired $V$.
This repairs both excluded carriers without assuming a standard
kernel whose endpoints would coincide or hit its hole.
\end{proof}

\begin{proposition}[Equal odd upper interval]\label{ok:equal-odd-offsets}
For every odd $L=2s+1\ge5$,
\[
 [0,3s-1]\ \subseteq\ \Om(\Sp{L,L}).
 \]
\end{proposition}
\begin{proof}
The basic equal-arm construction in Section~\ref{app:pair-offsets}
supplies $0\le k<L$. It remains to cover
$2s\le k\le3s-1$, an interval overlapping the basic one.
Set $V=k-2s+1$, so $1\le V\le s$.
Lemma~\ref{ok:variable-hole} supplies $K(r,s,s,V)$,
and extraction gives $k=s+s+V-1$.
All endpoint and label/color conditions are exactly those of
Lemma~\ref{ok:pair-extraction}; no additional path request is hidden.
\end{proof}

\subsubsection{Consecutive odd lengths}

\begin{lemma}[Consecutive-carrier kernels]\label{ok:consecutive-kernels}
For every $t\ge3$ and $1\le V\le t+1$, there is some $r$
and a kernel $K(r,t+1,t,V)$.
\end{lemma}
\begin{proof}
The fifteen initial words in Table~\ref{ok:consecutive-table}
cover every required $V$ at $t=3,4,5$.
For a row with parameters $t,V,r$, the missing label is $t+V$,
the missing color is $r+V-1$, and the two ordered endpoints are
$2t-r+2,t+2r-V+2$.
Listing the entries and subtracting successive pairs gives precisely
those inventories. Thus the table consists of full literal kernels,
rather than declarations of a finite search outcome.

For $t\ge6$ first note the following ordinary source. For
$T\ge1$, $T\ne2$, take the odd path on $[1,2T+1]$ with upper
endpoints $2T,T+2$, and raise upper labels by one.
Its unique forbidden request would satisfy $2(2T-1)=3T$,
namely $T=2$. Thus this gives
\begin{equation}
 K(1,T+1,T,1)\qquad(T\ge1,\ T\ne2).
 \label{ok:consecutive-ordinary-seed}
\end{equation}
The literal source
\[
 M=(6,1,8,2,4,3,7)=K(2,4,3,2)
 \]
will repair the missing $T=2$ case where it arises.

At $V=1$, use \eqref{ok:consecutive-ordinary-seed} with $T=t$.
At $V=2$, use carrier prefix of size $w=t-3\ge3$ on $M$.
Reflection preserves $U=t+1$ and sends $V$ to $t-V+3$;
therefore this also gives $V=t+1$.

For $3\le V\le t$ and $2V\ge t+3$, put
$T=t-V+1$ and $P=V-1$.
Then $T\ge1$ and $P\ge T+1$.
If $V\ne t-1$, then $T\ne2$ and diagonal growth of
\eqref{ok:consecutive-ordinary-seed} gives the desired kernel:
its new carrier is $T+P=t$, its other half is $t+1$,
and its hole parameter is $1+P=V$.
If $V=t-1$, diagonally grow $M$ by $P=t-3\ge3$ instead.
Its gap $T-S$ is $1$, so this growth is permitted and gives
$U=t+1$, carrier $t$, and $V=t-1$.
Finally, every remaining $3\le V\le t$ has $2V<t+3$;
reflect the already constructed upper-half kernel at
$V'=t-V+3$. It has $3\le V'\le t$ and $2V'\ge t+3$.
All hole parameters have now been covered.
\end{proof}

\begin{longtable}{@{}rrr@{\quad}l@{}}
\caption{Literal consecutive-carrier kernels.}
\label{ok:consecutive-table}\\
\toprule
$t$&$V$&$r$&Word\\ \midrule
\endfirsthead
\toprule $t$&$V$&$r$&Word\\ \midrule
\endhead
\bottomrule
\endfoot
$3$&$1$&$1$&$(7,2,8,1,5,3,6)$\\
$3$&$2$&$2$&$(6,1,8,2,4,3,7)$\\
$3$&$3$&$1$&$(7,2,8,1,5,3,4)$\\
$3$&$4$&$2$&$(6,2,8,1,4,3,5)$\\
$4$&$1$&$1$&$(9,1,10,3,6,4,8,2,7)$\\
$4$&$2$&$3$&$(7,4,5,3,8,2,9,1,10)$\\
$4$&$3$&$1$&$(9,3,8,1,10,2,6,4,5)$\\
$4$&$4$&$4$&$(6,4,5,2,7,3,9,1,10)$\\
$4$&$5$&$2$&$(8,1,10,2,7,3,6,4,5)$\\
$5$&$1$&$1$&$(11,2,12,1,9,3,10,5,7,4,8)$\\
$5$&$2$&$2$&$(10,2,11,1,12,5,6,4,8,3,9)$\\
$5$&$3$&$1$&$(11,2,12,1,9,4,10,3,7,5,6)$\\
$5$&$4$&$2$&$(10,4,8,1,12,2,11,3,6,5,7)$\\
$5$&$5$&$3$&$(9,1,12,2,11,5,6,4,7,3,8)$\\
$5$&$6$&$3$&$(9,4,8,2,12,1,10,3,6,5,7)$\\
\end{longtable}

\begin{proposition}[Consecutive odd upper interval]\label{ok:consecutive-offsets}
For odd $a\ge7$,
\[
 [a,(3a-1)/2]\ \subseteq\ \Om(\Sp{a,a+2}).
 \]
\end{proposition}
\begin{proof}
Write $a=2t+1$, $u=t+1$. The preceding lemma supplies every
$K(r,u,t,V)$ with $1\le V\le t+1=u$.
Extraction gives $k=2t+V$, ranging from $a$ to $3t+1$.
\end{proof}

\subsubsection{Two bands from the central and short-wedge kernels}

\begin{proposition}[Short-wedge pair band]\label{ok:short-pair-band}
Let $1\le t\le u$ and put
\[
 S=k-2u-2,\qquad\delta=u-t,\qquad V=k-t-u+1.
 \]
If
\[
 S\ge2,\quad t\ge2S+2,\quad
 0\le\delta\le S-1,\quad\delta\ne S-3,\quad1\le V\le t,
 \]
then $k\in\Om(\Sp{2t+1,2u+1})$.
\end{proposition}
\begin{proof}
Lemma~\ref{ok:unequal-wedge-kernel}, in its nonexceptional branch,
gives $K(S,t,u,S+\delta+3)$.
Here $S+\delta+3=k-t-u+1=V$.
Its ordinary odd seed has carrier $r=S+\delta+1$ and upper
first entry $2r-S+2$; the explicitly excluded request is
exactly $\delta=S-3$.
Its even carrier prefix has half-order $w=t-S-1\ge S+1$.
Thus all the source hypotheses have been checked, not inferred
from a general three-arm $\alpha$ theorem.
Apply extraction with carrier $u$ and other half $t$.
The two lengths are produced in reversed order, and its hypothesis
is precisely $1\le V\le t$. The offset is $u+t+V-1=k$.
\end{proof}

\begin{proposition}[Central pair band]\label{ok:central-pair-band}
Let $1\le t\le u$ and put $S=k-2t$, $V=k-t-u+1$.
If
\[
 1\le S\le t\le u\le2S-2,\qquad
 u>S,\quad u\ne2t-S,\quad1\le V\le u,
 \]
then $k\in\Om(\Sp{2t+1,2u+1})$.
\end{proposition}
\begin{proof}
The source of Lemma~\ref{ok:central-kernel} has
\[
 S_0=u-t,\quad T_0=u-S,\quad
 U_0=2u-S-t,\quad P=S+t-u.
 \]
The stated inequalities give $0\le S_0\le T_0$,
$T_0,U_0\ge1$, and $P\ge T_0-S_0+2$.
Its odd-path exception is exactly $u=2t-S$, which has been
excluded, and diagonal growth gives
$K(S,u,t,P+1)=K(S,u,t,V)$.
Extraction now gives the desired pair in its original order and
offset $t+u+V-1=k$. The extra $u>S$ avoids the equal zero-carrier
source; no AP or equal exception is silently included.
\end{proof}

\subsubsection{Asymmetric, strict one-layer, and reflected bands}

\begin{proposition}[Asymmetric pair bands]\label{ok:asymmetric-pair-bands}
Let $1\le t\le u$ and $V=k-t-u+1$.
Each of the following conditions suffices for
$k\in\Om(\Sp{2t+1,2u+1})$.
\begin{enumerate}
\item Put $s=k-2t$, $g=u-2s$, and $h=k-t-u$.
Require
\[
 g\ge1,\quad h\ge2,\quad s\ge g+h+1,\quad1\le V\le u.
 \]
\item Put $s=u-t+1$ and $g=2t-u-2$.
Require
\[
 k=t+u+2,\quad g\ge1,\quad s\ge g+3,\quad u\ge3.
 \]
\item Put
\[
 s=3u-k+3,\quad h=t+2u-k+1,\quad
 g=2k-4u-t-4.
 \]
Require
\[
 g\ge1,\quad h\ge2,\quad s\ge g+h+1,\quad1\le V\le t.
 \]
\end{enumerate}
\end{proposition}
\begin{proof}
For the first case substitute the displayed parameters into
\eqref{ok:transition-high}. They satisfy
\[
 2s+g=u,\qquad s+g+h=t,\qquad h+1=V.
 \]
Thus the kernel is $K(s,u,t,V)$ and extraction gives the claim.

For the second case substitute into \eqref{ok:transition-one}.
Again $2s+g=u$ and $s+g+1=t$, while $V=3$.
The required completion/extraction inequality is $3\le u$,
which is explicitly assumed. This is the strict $h=1$ source:
the limiting AP ray is not included by weakening $s\ge g+3$.

For the third case the source \eqref{ok:transition-high} has
\[
 T=s+g+h=u,\qquad U=2s+g=2u-t+2,\qquad W=h+1.
 \]
Reflect it using $R$. The resulting carrier remains $u$, its
other half is $2u-U+2=t$, and its hole parameter is
\[
 u-W+3=u-h+2=k-t-u+1=V.
 \]
The new shortest parameter is $u-s+1$; its value is already
realizable because it comes from a genuine reflected word.
Extraction is used in reversed arm order, and therefore needs
$1\le V\le t$, not merely $V\le u$.
This essential guard is retained in the statement.
\end{proof}

\subsection{Proof dependencies and computational verification}
\label{ok:verification}

The infinite arguments of this section consist of interval
partitions, the two classical prescribed-path interfaces, and the
controlled two-step recurrences. No finite search is used to infer
the existence of a path for arbitrarily large parameters.
Every literal initial word required by these arguments is printed
above. The negative $(3,3,3)$ profile is excluded by the degree and
color argument in Theorem~\ref{ok:all-odd-theorem}.

The proofs
of the graph constructions are the formulas and inventory arguments
in this section; they are not claimed to be formalized in a proof
assistant.
